% !PRINTED 48
\weberpage{48}
%<<CONT
by which the equations (16) are easily brought into the form:
\weq{17}{\begin{aligned}
\sqrt{x_1} &= -\dfrac{\sqrt{x_2}\sqrt{(1-x_3)(1-\varkappa^2x_3)} + \sqrt{x_3}\sqrt{(1-x_2)(1-\varkappa^2x_2)}}{1-\varkappa^2x_2x_3}, \\
\sqrt{1-x_1} &= \dfrac{\sqrt{1-x_2}\sqrt{1-x_3} - \sqrt{x_2x_3}\sqrt{(1-\varkappa^2x_2)(1-\varkappa^2x_3)}}{1-\varkappa^2x_2x_3}, \\
\sqrt{1-\varkappa^2x_1} &= \dfrac{\sqrt{1-\varkappa^2x_2}\sqrt{1-\varkappa^2x_3} - \varkappa^2\sqrt{x_2x_3}\sqrt{(1-x_2)(1-x_3)}}{1-\varkappa^2x_2x_3},
\end{aligned}}
and our consideration teaches that, if $x_1$ and the roots
$\sqrt{x_1}, \sqrt{1-x_2}, \sqrt{1-\varkappa^2x_1}$ % [sic] source lists \sqrt{1-x_2}; equations (17) and this x_1 list require \sqrt{1-x_1}.
are determined by these equations as functions
of $x_2, x_3$, the differential equation (1) is identically
satisfied.

\begin{center}\rule{2cm}{0.4pt}\end{center}

In order to obtain the addition theorem in the simplest possible
 shape for the Weierstrass normal form as well, we set
\[ f(x) = 4x^3 - g_2x - g_3, \]
and let $x_4$ become infinitely great in equation (11) (after division
 by $x_4^2$). There then results the integral of the
differential equation (1)
\[ \dfrac{dx_1}{\sqrt{f(x_1)}} + \dfrac{dx_2}{\sqrt{f(x_2)}} + \dfrac{dx_3}{\sqrt{f(x_3)}} = 0 \]
in the form
\weq{18}{\begin{vmatrix} 1, & x_1, & \sqrt{f(x_1)} \\ 1, & x_2, & \sqrt{f(x_2)} \\ 1, & x_3, & \sqrt{f(x_3)} \end{vmatrix} = 0,}
or
\weq{19}{a + bx_1 + c\sqrt{f(x_1)} = 0,}
if
\weq{20}{\begin{aligned}
a &= x_2\sqrt{f(x_3)} - x_3\sqrt{f(x_2)} \\
b &= \sqrt{f(x_2)} - \sqrt{f(x_3)} \\
c &= x_3 - x_2.
\end{aligned}}
Then $x_1, x_2, x_3$ become the roots of the cubic
equation
\[ (a+bx)^2 - c^2f(x) = 0, \]
and if we here set the coefficient of $x^2$, divided by
that of $x^3$, equal to the negative sum of the roots, we
obtain
%CONT>>
